\documentclass[10pt,a4paper]{amsart}
\usepackage[margin=19mm]{geometry}
\usepackage{amsmath,amssymb,mathtools}
\usepackage[scr=rsfs,cal=euler]{mathalfa}
\usepackage{xfrac}
\usepackage[unicode,hidelinks,bookmarks=false]{hyperref}
\newcommand{\ie}{i.e.\ }
\newcommand{\cf}{cf.\ }
\newcommand{\resp}{respectively\ }
\newcommand{\Subsection}{\subsection}
\providecommand{\nobreakdash}{\nobreak\hbox{-}}
\setlength{\emergencystretch}{2em}
\multlinegap0pt
\usepackage[all]{xy}
\usepackage{float}
\usepackage{mathtools}
\usepackage{amssymb}
\usepackage{enumerate}
\usepackage{xcolor}
\usepackage{tikz}
\usepackage{bbm}
\usepackage[shortlabels]{enumitem}
\newtheorem{lemma}{Lemma}
\title{Irreducible objects in the Gaiotto category at roots of unity}
\author{Aleksandr Popkovich}
\date{Source: arXiv:2602.08264v1 (CC BY 4.0)}
\begin{document}
\maketitle

\noindent\emph{Source and licence.} Irreducible objects in the Gaiotto category at roots of unity. Version arXiv:2602.08264v1 (CC BY 4.0), distributed by arXiv under the Creative Commons Attribution 4.0 International licence, http://creativecommons.org/licenses/by/4.0/, as recorded in the arXiv licence metadata for this exact version and shown on the abstract page https://arxiv.org/abs/2602.08264v1. Full licence notice: \texttt{licenses/arxiv-2602.08264-CC-BY-4.0.txt}.\par\smallskip

\noindent\textit{Editorial scope.} Counted unit: the article's lemma lemma (source0.tex lines 439--448). The article's complete original proof of it is delivered with it (source0.tex lines 450--479, one \texttt{proof} environment of 30 lines). No mathematics is written, repaired or re-derived: the statement and the proof are the article's own lines, in the article's own notation, and no retained line is substituted or re-flowed.  No further source-local context is needed: every cross-reference of every retained line resolves inside the two delivered spans.  Nothing was omitted from any retained span, so the package carries no omission record.  No retained line contains a citation, so the wrapper carries no bibliography and no bibliography entry.  No retained line contains a figure environment, an image-inclusion command or drawing code, so no image is delivered and the package declares no asset dependency.  Editorial formatting only, in addition: an AMS \texttt{amsart} preamble that restates the article's own declarations and macros used by the retained lines, namely the environments \texttt{lemma} on separate counters, together with the standard packages those lines need.  The retained environments print in this document as Lemma 1 in order of appearance, whereas the article prints the same items with the numbering of its own full document.  The complete native source of the article as distributed in the arXiv e-print for this exact version is bundled unchanged as \texttt{sources/194237/source0.tex}.
\begin{lemma}
Denote by $M \subset A$ the subset consisting of pairs $(\lambda, \theta)$ that satisfy the following  conditions:
\begin{enumerate}
    \item $\lambda_1 \ge \lambda_2 \ge ... \ge \lambda_M$ and $\theta_1\ge \theta_2 ... \ge \theta_{M+1}$
    \item if $\lambda_{i-1} =\lambda_i$ then $\lambda_i+\theta_i=0 \text{ mod }p$, and
    \item if $\theta_{i} =\theta_{i+1}$ then $\lambda_i+\theta_i=0 \text{ mod }p$.
\end{enumerate}

    Then we have an equality of sets $S(A)=M \subset \mathbb{Z}^M\times\mathbb{Z}^{M+1}$.
\end{lemma}

\begin{proof}
    First, let us show that $S(A)\subset M$. We will use the fact that the linear ordering of $\Phi^+_{st}\backslash \Phi^+_m$ can be chosen in any way which is consistent with the partial order on this set without affecting the result of the algorithm.
    \begin{itemize}
        \item To show that any element $(\tilde\lambda,\tilde\theta)=S\left( (\lambda,\theta)\right)$ in the image of $S$ has $\tilde\lambda_{i} \ge \tilde\lambda_{i+1}$ for all $i$  and the condition (i) of the Lemma is satisfied, we use Version 1 of the algorithm.

        First we show that $\tilde\lambda_{i} \ge \tilde\lambda_{i+1}$ for all $i$, that is $\tilde\lambda$ is a non-increasing sequence of integers. In fact, a stronger statement is true (and we will need it swiftly): the property of being non-increasing is preserved on every step of the algorithm. Indeed, for the sake of contradiction suppose that on $k$-th step of the algorithm this property is broken for the first time, that is $\lambda^{(k)}$ was non-increasing but $\lambda^{(k+1)}$ is not. Since only the $i_k$-th entry of $\lambda$ could be affected on this step, it must have been changed (decreased) and we must have $\lambda^{(k+1)}_{i_k}<\lambda^{(k+1)}_{i_k+1}$. Since every time Action 2 is taken the entry is only decreased by one, and on the previous iteration $\lambda^{(k)}$ was non-increasing, we must have had $\lambda^{(k)}_{i_k}<\lambda^{(k)}_{i_k+1}$, and on a previous $(k-1)$-th step with $i_{k-1}=i_k+1$ (since we are using the dirst ordering) Action 1 must have been taken. But then it also should have been taken on the $k$-th step. A contradiction.

         Now to prove that condidtion (i) is always satisfied for $(\tilde\lambda,\tilde\theta)$, suppose that we have $\tilde\lambda_{i-1} =\tilde\lambda_{i}$ for some $i$. Let $k$ be the last step on which $i_k=i$ so that $\lambda^{(k+1)}_{i}=\tilde\lambda_{i}$. Since this is the last step on which $i_k=i$ and we are using the ordering $<_1$ the following must be true: we must have $j_k=i_k$, and all of the remaining ``m"-th steps have $i_m>i_k$ and $j_m>j_k$. It follows that $\lambda_{i_k-1}^{(k+1)}= \tilde\lambda_{i_k-1}$, $\lambda_{i_k}^{(k+1)}= \tilde\lambda_{i_k}$ and $\theta_{i_k}^{(k+1)}= \tilde\theta_{i_k}$. So due to our assumption we must have $\lambda_{i-1}^{(k+1)} =\lambda_{i}^{(k+1)}$. This, together with the fact (proven above) that $\lambda^{(k)}$ was non-increasing implies that on $k$-th step Action 1 was taken, so $\lambda_{i_k}^{(k+1)}= \theta_{i_k}^{(k+1)}$ and hence $\tilde\lambda_{i_k}^{(k+1)}= \tilde\theta_{i_k}^{(k+1)}$.

        
        \item Analogous reasoning using Version 2 of the algorithm shows that any element $(\tilde\lambda,\tilde\theta)=S\left( (\lambda,\theta)\right)$ in the image of $S$ has $\tilde\theta_{i} \ge \tilde\theta_{i+1}$ for all $i$  and the condition (ii) of the Lemma is satisfied.

        \item Finally, let us show that any pair $(\tilde\lambda,\tilde\theta) \in M$ can be obtained as $S((\lambda,\theta))$ for some $(\lambda,\theta) \in A$. It turns out one can describe an inverse map $S^{-1}:M\to A$.

        For this note that whichever of the two Actions of the algorithm is take on step $k$, we have $\lambda^{(k)}_{i_k}+\theta^{(k)}_{j_k} \equiv \lambda^{(k+1)}_{i_k}+\theta^{(k+1)}_{j_k} \text{ mod }p$; so at every moment we "remember" what Action we took on the previous step. Therefore, we can "invert" the algorithm by going through the set $\Phi^+_{st}\backslash \Phi^+_m$ in the order inverse to either $<_1$ or $<_2$ and performing the transformations inverse to the ones we did earlier.

        This defines a map $S^{-1}:M \to \mathbb{Z}^M\times\mathbb{Z}^{M+1}$; we need to show that $S^{-1}(M)\subset A$, i.e. that for any element $(\tilde\lambda,\tilde\theta)\in M$ and $(\lambda,\theta):=S^{-1}((\tilde\lambda,\tilde\theta))$ both $\lambda$ and $\theta$ are non-decreasing.

        Let us prove that $\lambda$ is non-decreasing.  For this we will use the inverse to the Version 1 of the Serganova algorithm and the condition (i) of the Lemma. Since $\tilde\lambda$ is already non-decreasing, we just need to show that this property is preserved on each step of the procedure. 

        Let $k$-th step of the reversed algorithm be the first one on which the property of being non-increasing is broken. Call the input of this step $(\lambda^{(in)},\theta^{(in)})$ and the output $(\lambda^{(out)},\theta^{(out)})$. The property being broken means that there are indeces $i$ and $j$ (indeces of entries being considered on this step) such that  $\lambda^{(in)}_{i-1}=\lambda^{(in)}_{i}$, but $\lambda^{(out)}_i= \lambda^{(out)}_{i-1}+1$ and $\theta^{(out)}_j = \theta^{(in)}_j-1$, that is on this step the inverse of Action 2 was performed, increasing the $i$-th value of $\lambda$ and decreasing the $j$-th values of $\theta$. Now consider two cases: either $i=j$ or $i>j$. 

        If $i>j$ then the above could not happen: for the inverse of Action 2 to be taken we must have $\lambda^{(in)} + \theta^{(in)} \ne 0\text{ mod }p$, but to have $\lambda^{(in)}_{i-1}=\lambda^{(in)}_{i}$ Action 1 must have been taken on the previous step  - one easily sees a contradiction.

        If $i=j$  then again the inverse to Action 2 could not have been performed on this step: due to how reverse to order $<_1$ works, we must have $\lambda^{(in)}_{i-1}=\tilde\lambda_{i-1} $, $\lambda^{(in)}_{i}= \tilde\lambda_{i}$ and $\theta^{(in)}_{i}= \tilde\theta_{i}$, and we know that any element $(\tilde\lambda,\tilde\theta)\in S(A)$ satisfies condition (i) of the Lemma, so from $\lambda^{(in)}_{i-1}=\lambda^{(in)}_{i}$ follows $\lambda^{(in)} + \theta^{(in)} = 0\text{ mod }p$ hence the inverse to Action 2 was not taken.

        Proof that $\theta$ is non-decreasing is analogues, using the inverse to Version 2 of the algorithm and condition (ii) of the lemma.

    \end{itemize}
\end{proof}

\end{document}
